\documentclass[conference]{IEEEtran}
\IEEEoverridecommandlockouts

\usepackage[utf8]{inputenc}
\usepackage[T1]{fontenc}
\usepackage{cite}
\usepackage{amsmath,amssymb,amsfonts}
\usepackage{graphicx}
\usepackage{xcolor}
\usepackage{booktabs}
\usepackage{textcomp}
\usepackage{hyperref}

\hypersetup{colorlinks=true,linkcolor=blue,citecolor=blue,urlcolor=blue}

\def\BibTeX{{\rm B\kern-.05em{\sc i\kern-.025em b}\kern-.08em
    T\kern-.1667em\lower.7ex\hbox{E}\kern-.125emX}}

\begin{document}

\title{T\"urkResearcher: A LangGraph-based Multi-Agent\\Research Assistant Grounded in 633K Turkish Theses}

\author{
\IEEEauthorblockN{Hakan Sabuni\c{s}}
\IEEEauthorblockA{Department of Computer Engineering\\
Istanbul Medipol University\\
Email: hakansabunis@gmail.com}
}

\maketitle

\begin{abstract}
Despite the rapid emergence of English-language research assistants such as
Elicit and Consensus.app, no comparable system exists for Turkish academic
literature, leaving researchers without grounded, citation-supported answers
in their native language. We introduce \textit{T\"urkResearcher}, a multi-agent
research assistant built on LangChain and LangGraph, which decomposes Turkish
academic queries, retrieves evidence from a corpus of 633{,}998 Turkish thesis
abstracts (the largest publicly indexed corpus to date for the domain),
synthesises findings, and produces IEEE-style citations. The system uses a
hybrid architecture: a multilingual sentence-transformer embedder for offline
retrieval and the DeepSeek API for reasoning. We additionally harvest
hundreds of thousands of journal articles from DergiPark via OAI-PMH to
extend coverage. We evaluate the system on a 30-question Turkish-language
benchmark spanning ten subject areas using LLM-as-judge metrics for citation
accuracy, faithfulness, sub-question coverage, and overall academic quality.
On a 30-question Turkish benchmark spanning ten subject categories, the
633K-thesis configuration of the agent achieves a median citation
accuracy of 0.70 and a median faithfulness of 0.63, with strong
performance in domains where Turkish theses are densely indexed
(health, education, engineering) and a clear failure mode in computer
science where Turkish researchers publish predominantly in English. We
additionally evaluate the 740K-record corpus produced by adding
106{,}641 DergiPark journal abstracts and report a counter-intuitive
result: naive expansion improves precisely the under-served categories
(CS +0.05, business +0.20) but \emph{regresses} the well-covered ones
(health $-0.27$, engineering $-0.20$) due to abstract-length distribution
shift and citation inflation. The full pipeline, the 30-question
benchmark, and the prebuilt 16\,GB Chroma index are released openly on
the Hugging Face Hub for reproducibility.
\end{abstract}

\begin{IEEEkeywords}
Retrieval-Augmented Generation, Multi-Agent Systems, Turkish NLP, LangChain,
LangGraph, Academic Search, LLM-as-Judge.
\end{IEEEkeywords}

\section{Introduction}

\IEEEPARstart{R}{esearch} assistants powered by retrieval-augmented generation
(RAG) and large language models (LLMs) have rapidly become an essential layer
of modern academic workflows. English-language tools such as Elicit, scite.ai
and Consensus.app help researchers locate evidence, distil findings, and trace
arguments to primary sources. None of these systems serve the Turkish-language
academic ecosystem. Turkish researchers, despite the country having more than
600 active scholarly journals on DergiPark and tens of thousands of theses
released annually through YÖK Ulusal Tez Merkezi, must currently choose
between general-purpose chatbots that hallucinate or English-only tools that
ignore the Turkish corpus entirely.

This work fills that gap with \textit{T\"urkResearcher}, an open multi-agent
academic research assistant for Turkish. Given a research question in Turkish,
the system decomposes it into sub-questions, retrieves evidence multi-querily
from a 633{,}998-record corpus of Turkish thesis abstracts, synthesises the
findings, runs an explicit critic step, and finally drafts a Turkish academic
answer with IEEE-style citations grounded in retrieved sources.

\textbf{Contributions.}
\begin{itemize}
\item A \emph{novel application} of multi-agent RAG to the under-served Turkish
academic-search niche, with no commercial or open-source equivalent at the
time of writing.
\item A reproducible \emph{data pipeline} built atop the
\texttt{umutertugrul/turkish-academic-theses-dataset} corpus, applying
quality filters and producing a \emph{title-aware}, cosine-distance Chroma
index that we release publicly as
\texttt{hakansabunis/tr-academic-research-agent-index}.
\item An \emph{OAI-PMH harvester} for DergiPark that extends the corpus by
several hundred thousand journal articles.
\item A \emph{LangGraph state machine} that integrates Planner, Retriever,
Synthesiser, Critic and Writer agents, and demonstrates conditional looping
when retrieval coverage is insufficient.
\item A \emph{30-question Turkish benchmark} and an LLM-as-judge evaluation
framework reporting citation accuracy, faithfulness, multi-hop coverage, and
holistic academic quality.
\end{itemize}

\section{Related Work}

\textbf{RAG and research assistants.} Retrieval-augmented generation
\cite{lewis2020rag} has become the dominant paradigm for grounding LLM
output in external evidence. Commercial systems such as Elicit and
Consensus.app demonstrate the demand for evidence-traceable academic search,
but they target English-only corpora. Open-source efforts such as
\texttt{LangChain} and \texttt{LangGraph} \cite{langchain2024,langgraph2024}
provide composable RAG and multi-agent abstractions that we build upon.

\textbf{Turkish NLP.} Recent work on Turkish NLP has produced strong
language-model baselines (BERTurk \cite{schweter2020berturk}, multilingual
sentence transformers \cite{reimers2019sbert,reimers2020multilingualsbert})
and several monolingual datasets, but coverage of academic search and
domain-specific RAG remains thin. The
\texttt{umutertugrul/turkish-academic-theses-dataset} corpus
\cite{umut2024theses} provides $\sim$650K bilingual abstracts that we filter
and index in this work; previous efforts such as TR-Akademik~\cite{trakad2025}
indexed only a $\sim$48K subset using a leaner schema and an L2 distance
metric, which we improve upon.

\textbf{Multi-agent LLM systems.} Recent work has shown that decomposing a
task across specialised agents (e.g., AutoGen \cite{wu2023autogen}) yields
better results than monolithic prompting, particularly when retrieval and
critique steps are interleaved. We instantiate a five-agent pipeline with
explicit conditional looping inspired by these systems.

\textbf{Evaluation of RAG.} Frameworks such as RAGAS \cite{es2024ragas} and
RAG\,Bench \cite{lyu2024rag} measure faithfulness, answer relevance and
context precision. We adopt an LLM-as-judge approach \cite{zheng2023llmjudge}
adapted to Turkish, scoring citation accuracy, faithfulness, multi-hop
coverage, and overall academic quality.

\section{Methodology}

\subsection{System overview}

T\"urkResearcher is a directed graph of five agents orchestrated by
LangGraph, with a single conditional edge from \texttt{Critic} back to
\texttt{Retriever} (Fig.~\ref{fig:arch}). All four LLM-using agents share the
same DeepSeek backend through LangChain's \texttt{ChatOpenAI} adapter,
configured to use \texttt{function\_calling} for structured output, since
DeepSeek does not yet expose the OpenAI \texttt{json\_schema} response
format.

\begin{figure}[t]
\centering
% \includegraphics[width=\linewidth]{figures/architecture.pdf}
\fbox{\parbox{0.92\linewidth}{\centering
\textbf{Question (TR)} $\to$
\textbf{Planner} (3--5 sub-Qs) $\to$
\textbf{Retriever} (multi-query Chroma) $\to$
\textbf{Synthesiser} (cluster findings) $\to$
\textbf{Critic} (coverage check) $\to^{\mathrm{loop}}$
\textbf{Writer} (IEEE-cited answer)
}}
\caption{T\"urkResearcher LangGraph state machine. The Critic node may route
back to Retriever up to twice if coverage is insufficient.}
\label{fig:arch}
\end{figure}

\subsection{Data pipeline}

We start from \texttt{umutertugrul/turkish-academic-theses-dataset}, a
$\sim$650K-row bilingual thesis-metadata dataset under CC-BY-4.0, and apply
the following filters:
(i) \texttt{abstract\_tr} non-empty;
(ii) \texttt{abstract\_tr} $\geq 50$ tokens;
(iii) \texttt{title\_tr} non-empty;
(iv) deduplication by \texttt{tez\_no}.
This yields $N{=}633{,}998$ records spanning 1966--2024 across more than
50 distinct subject areas (Education, Business, Agriculture, Law, Medicine,
Engineering, etc.; the most populous subjects are Education and Training
($n{=}5126$) and Business Administration ($n{=}2048$)).

Each record is encoded as the concatenation
\texttt{title\_tr}~$\Vert$~\texttt{abstract\_tr}, providing a stronger
lexical signal than the abstract alone, and embedded with
\texttt{paraphrase-multilingual-mpnet-base-v2} (768-dim). We persist
embeddings into Chroma using cosine distance (\texttt{hnsw:space=cosine}).

\subsection{Multi-agent pipeline}

\textbf{Planner.} Decomposes the user's question into 3--5 well-typed
Turkish sub-questions emphasising different conceptual angles (method,
finding, limitation, comparison, application). Output is a Pydantic
\texttt{Plan} object enforced via function calling.

\textbf{Retriever.} Runs the original question and every sub-question
against the Chroma collection (top-$k{=}6$ each), deduplicates by
\texttt{tez\_no} keeping the highest score, and returns up to 24 chunks
sorted by similarity. Cosine distance is converted into a
$[0,1]$ similarity score.

\textbf{Synthesiser.} Receives the retrieved chunks and produces a
\texttt{Synthesis} object whose \texttt{findings} are short claims with
explicit citations to chunk-level \texttt{tez\_no}s, plus a
\texttt{contradictions} list when the corpus disagrees with itself.

\textbf{Critic.} Evaluates whether the synthesis covers all sub-questions,
emitting Boolean \texttt{coverage\_ok}, a list of missing aspects, and
concrete Turkish \texttt{requery\_terms} for the next retrieval pass. If
\texttt{coverage\_ok} is false and the iteration count is below
\texttt{MAX\_CRITIC\_LOOPS}~$=2$, control returns to the Retriever; otherwise
flow proceeds to the Writer.

\textbf{Writer.} Produces a final Turkish academic answer in Markdown,
sentence-level citations \texttt{[n]} matching the chunks, and IEEE-style
references that include author, title, location, and the public YÖK PDF
URL.

\subsection{DergiPark extension}

To extend coverage beyond theses, we implement an OAI-PMH harvester
(\texttt{scripts/05\_harvest\_dergipark.py}) that streams Dublin Core records
from
\url{https://dergipark.org.tr/api/public/oai/}, persists JSONL incrementally
with resumption-token state for crash safety, and respects 429 backoff. At
the time of writing this corpus is being embedded and merged into the same
Chroma collection.

\section{Experiments}

\subsection{Setup}

\textbf{Corpus index.} Final 633{,}998-document Chroma index, built on a
Colab T4/A100 GPU with mpnet-base-v2 over $\sim$6 hours, occupying 13.4 GB
on disk and 14.8 GB on Hugging Face Hub.

\textbf{Hardware.} Local agent inference on a Windows 11 laptop with an RTX
3050 Ti (4 GB) acting as embedder host (CPU also sufficient); LLM calls
issued to the DeepSeek API (\texttt{deepseek-chat}). No GPU is required at
query time once the index is downloaded.

\textbf{Benchmark.} A 30-question Turkish-language eval set,
\texttt{data/eval/questions.json}, with expected subject areas and minimum
citation thresholds. Categories cover computer science, education, health,
engineering, social sciences, business, agriculture, veterinary, linguistics,
law, physics, chemistry, biology, tourism, sports and cross-domain queries.

\subsection{Metrics}

For each question, we compute four metrics with DeepSeek as judge:
\begin{itemize}
\item \textbf{Citation accuracy} ($\in [0,1]$): the fraction of in-text
\texttt{[n]} citations whose underlying chunk supports the surrounding claim.
\item \textbf{Faithfulness} ($\in [0,1]$): how grounded the answer is in
retrieved chunks (1.0 = no hallucination).
\item \textbf{Coverage} ($\in [0,1]$): fraction of planner sub-questions
substantively addressed in the final answer.
\item \textbf{Holistic score} ($\in \{1,\dots,5\}$): overall academic quality
(structure, register, reasoning).
\end{itemize}
We additionally report end-to-end latency, number of retrieved chunks,
maximum chunk similarity, and the number of citations.

\subsection{Results}

We report headline numbers for two corpus configurations:
the original 633{,}998-thesis index, and an extended 740{,}639-record
index produced by appending 106{,}641 Turkish journal articles harvested
from DergiPark via OAI-PMH (\S\ref{sec:dergipark}). All 30 benchmark
questions completed without agent failures in both runs. Each question
retrieved on average 30 chunks (min 24, max 40), with a top similarity
score of 0.93 (cosine) on the 633K corpus, indicating that the retriever
consistently surfaces well-matched evidence in both regimes.

Table~\ref{tab:overall} reports the headline LLM-as-judge metrics, and
Table~\ref{tab:bycat} breaks them down by question category.

\begin{table}[t]
\centering
\caption{Headline metrics: 633K theses vs 740K (+ DergiPark) corpus, mean
over n=30 questions. $\Delta$ is 740K minus 633K.}
\label{tab:overall}
\begin{tabular}{lccc}
\toprule
Metric & 633K & 740K & $\Delta$ \\
\midrule
Citation accuracy & 0.60 & 0.51 & $-0.10$ \\
Faithfulness      & 0.59 & 0.49 & $-0.10$ \\
Coverage          & 0.49 & 0.47 & $-0.03$ \\
Holistic (1--5)   & 2.63 & 2.40 & $-0.23$ \\
\#Citations       & 30.1 & 32.8 & $+2.7$ \\
Latency (s)       & 105.8& 111.4& $+5.6$ \\
\bottomrule
\end{tabular}
\end{table}

\begin{table}[t]
\centering
\caption{Per-category citation accuracy, sorted by 740K--633K delta.
Bold rows are categories that improved with corpus expansion;
others regressed. $n=1$ categories are omitted (insufficient sample).}
\label{tab:bycat}
\begin{tabular}{lcccc}
\toprule
Category & n & 633K & 740K & $\Delta$ \\
\midrule
\textbf{business}        & 2 & 0.47 & 0.68 & $+0.20$ \\
\textbf{computer\_science}& 4 & 0.21 & 0.26 & $+0.05$ \\
\midrule
social\_sciences         & 2 & 0.78 & 0.70 & $-0.08$ \\
education                & 3 & 0.75 & 0.63 & $-0.12$ \\
multi\_domain            & 2 & 0.85 & 0.70 & $-0.15$ \\
engineering              & 3 & 0.73 & 0.53 & $-0.20$ \\
law                      & 2 & 0.62 & 0.40 & $-0.22$ \\
health                   & 3 & 0.80 & 0.53 & $-0.27$ \\
\bottomrule
\end{tabular}
\end{table}

\textbf{The corpus-expansion paradox.} Adding 106K journal articles
\emph{regressed} mean citation accuracy from 0.60 to 0.51, and mean
faithfulness from 0.59 to 0.49 (Table~\ref{tab:overall}), while the
mean number of citations per answer increased from 30.1 to 32.8.
Per-category (Table~\ref{tab:bycat}) the picture is more nuanced:
the two original failure modes---\emph{computer\_science} and
\emph{business}---improved by +0.05 and +0.20 respectively, exactly
the categories where Turkish theses were sparse. But categories with
already-strong coverage (\emph{health}, \emph{engineering}, \emph{law})
regressed by 0.20--0.27. We discuss this in Section~\ref{sec:cs-failure}.

\section{Discussion and Error Analysis}

\textbf{Strengths.} Adding the title to the embedded document and switching
from Euclidean to cosine distance materially improves retrieval relevance
compared with the prior 48K-record TR-Akademik index. Average top-result
similarity is 0.93 (cosine) across the 30-question benchmark, and the
agent consistently produces 24--40 IEEE citations with public YÖK PDF
URLs. In broad domains where the corpus is well represented---health,
education, engineering, social sciences and multi-domain queries---the
median citation accuracy is 0.70+ and the holistic LLM-judge score is 3+,
which we consider a strong baseline for an unsupervised, zero-shot RAG
system over a 633K-document Turkish corpus.

\subsection{The Computer-Science Failure Mode}
\label{sec:cs-failure}

The single largest source of error in the benchmark is the computer-science
category (n=4): mean citation accuracy of 0.21 and a holistic score of
1.5/5, dragging the global mean from 0.70 (median) to 0.60 (mean). Manual
inspection of the failed runs reveals two compounding causes:

\begin{enumerate}
\item \textbf{Linguistic concentration of ``Türkçe'' theses.} A query such
as ``Türkçe doğal dil işlemede son yaklaşımlar'' over-retrieves theses on
the \emph{Turkish language itself}---phonology, dialectology, language
education---because the lexical overlap of the word ``Türkçe'' is
dominant in non-CS theses about the language. The cosine retriever, even
title-aware, cannot reliably distinguish ``Turkish (NLP)'' from
``Turkish (linguistics).''

\item \textbf{Corpus skew.} Computer-science theses in the YÖK collection
are sparser than expected, partly because Turkish CS researchers
disproportionately publish in English at international venues
(SIU, NeurIPS, ACL etc.), and the YÖK corpus is filtered to Turkish
abstracts only.
\end{enumerate}

Both problems would be ameliorated by (i) tighter retrieval through a
domain-tuned embedder (Section~\ref{sec:future}), and (ii) by routing
CS queries to OpenAlex, which we measured to contain
$\sim$1.16M Turkish-language papers spanning IEEE Xplore, SIU, Springer
and Elsevier venues. The live-search node already implements this
routing; in future work the Critic should learn to invoke it more
aggressively for CS-flagged queries.

\subsection{Why naive corpus expansion hurts well-covered domains}

Adding the DergiPark articles improved the under-served domains (CS, business)
but \emph{regressed} categories that already had dense thesis coverage
(health, engineering, law). Three concrete contributing factors emerge from
inspecting per-question outputs:

\begin{enumerate}
\item \textbf{Abstract length distribution shift.} Thesis abstracts in our
corpus have a median length of $\sim$1{,}600 characters, with the
title-aware document running 2{,}032 characters at the median. DergiPark
journal abstracts are typically 300--800 characters: less detail per
chunk, less surface area for the writer agent to ground a claim against.

\item \textbf{Citation inflation.} The writer produces 2.7 \emph{more}
citations per answer on average (32.8 vs 30.1). Each extra citation is
either weakly grounded (lowering citation accuracy) or weakly synthesised
(lowering faithfulness), even if the underlying chunk is on-topic.

\item \textbf{Source-mixing without source-aware writing.} Theses and
journal articles serve different rhetorical roles---a thesis abstract
states a single coherent contribution, a journal article often a tighter
empirical claim. Our writer prompt does not currently distinguish them,
so it sometimes glues a thesis-style overview claim to a narrow journal
finding, producing a fluent but partially unsupported sentence.
\end{enumerate}

The take-away is not that DergiPark integration was a mistake---it is the
right call for under-served domains and we keep it---but that
\emph{naive corpus expansion is not a free win}. The next iteration of the
system should make retrieval and writing source-aware: when local thesis
coverage is dense (e.g.\ health), prefer thesis chunks; when it is sparse
(e.g.\ CS), pull in journal articles and live OpenAlex/Semantic Scholar
results, with the writer prompt adjusted accordingly. We sketch this in
Section~\ref{sec:future}.

\textbf{Latency.} End-to-end runs are dominated by sequential DeepSeek calls
(Plan, Synth, Critic, Write), with a benchmark mean of 105.8\,s and a
median of 111.6\,s per question. The fastest run (q03) completed in 46\,s
and the slowest (q24) in 146.5\,s. Latency could be reduced by parallelising
sub-question retrievals (currently sequential) and by caching the Planner
output for semantically similar questions.

\textbf{Limitations.}
(i)~The corpus is restricted to thesis abstracts (full text is not yet
included); DergiPark integration partially addresses this with journal
articles. (ii)~The current evaluation uses a single LLM judge; cross-judge
agreement should be measured in future work. (iii)~The pipeline assumes
\emph{static} corpus knowledge between rebuilds; live API tools (DergiPark
REST, OpenAlex, Semantic Scholar) are flagged as future work for real-time
coverage of newly published research.

\subsection{Future Work: A Fine-Tuned Turkish Academic LLM}
\label{sec:future}

The current system relies on a general-purpose LLM (DeepSeek) operating in a
zero-shot RAG setting. The same data we have already curated--633K thesis
abstracts plus a parallel journal-article harvest--is, however, sufficient
to train a domain-specific Turkish academic model in three concrete steps.

\textbf{Stage 1 (custom embedder).} A SimCSE or
MultipleNegativesRankingLoss fine-tune of
\texttt{paraphrase-multilingual-mpnet-base-v2} on the 633K thesis title /
abstract pairs would specialise the embedding space to Turkish academic
register. We expect a 15--25\% relative improvement in retrieval relevance,
with a wall-clock cost of $\sim$3--5 hours on a Colab T4.

\textbf{Stage 2 (instruction-tuned LLM).} We can synthesise
$\sim$100K--200K \emph{grounded} question--answer pairs by prompting
DeepSeek (or GPT-4o) with each thesis abstract as context, then fine-tune
a Turkish base model (e.g.\ \texttt{Trendyol-LLM-7B}, \texttt{Cosmos-LLaMa},
or a multilingual base such as \texttt{Llama-3.2-3B}) using QLoRA. The
target dataset format is \texttt{\{question, retrieved\_chunks, grounded\_answer\}},
mirroring the agent's runtime pipeline. Training cost on an A100 is in the
order of 6--12 hours; inference quantises down to 4-bit and fits a 4 GB
laptop GPU.

\textbf{Stage 3 (preference alignment).} The 30 evaluation runs already
produce, via the LLM-as-judge pipeline, paired
$(\text{answer}, \text{judgment})$ data. Combining successful (high
citation accuracy + faithfulness) and unsuccessful answers yields preference
pairs suitable for Direct Preference Optimisation (DPO). This would teach
the model to prefer grounded, well-cited responses without further
human annotation.

The ultimate artefact would be \texttt{T\"urkResearcher-7B-instruct}: an
openly released Turkish academic assistant model that runs the full
multi-agent loop offline. Beyond the present project, this trajectory aligns
naturally with TÜBİTAK and YÖK initiatives on Turkish-language scientific
infrastructure, and would constitute a contribution well-suited to a
follow-up MSc/PhD thesis.

\section{Conclusion}

We presented \textit{T\"urkResearcher}, an open multi-agent academic research
assistant for Turkish, grounded in 633{,}998 thesis abstracts and a parallel
DergiPark journal harvest. The system targets a clear, previously unfilled
niche, integrates Planner/Retriever/Synthesiser/Critic/Writer agents in a
LangGraph state machine, and emits IEEE-style citations traceable to public
YÖK PDFs. The complete pipeline, including the 14.8\,GB prebuilt index, is
released on the Hugging Face Hub for reproducibility. Future work targets
live source aggregation, full-text retrieval, and cross-judge evaluation.

\bibliographystyle{IEEEtran}
\bibliography{refs}

\end{document}
